% Reconnaissances - Bourses de doctorat en recherche FRQ 2027-2028 (B2)
% Compiler avec XeLaTeX ou LuaLaTeX pour obtenir le vrai Times New Roman exigé par le FRQ.
% Avec pdfLaTeX, le document utilise un clone de Times (newtxtext).
\documentclass[12pt,letterpaper]{article}

\usepackage{iftex}
\ifPDFTeX
  \usepackage[T1]{fontenc}
  \usepackage[utf8]{inputenc}
  \IfFileExists{newtxtext.sty}{\usepackage{newtxtext}}{\usepackage{mathptmx}}
\else
  \usepackage{fontspec}
  \IfFontExistsTF{Times New Roman}{\setmainfont{Times New Roman}}{\setmainfont{TeX Gyre Termes}}
\fi
\usepackage[french]{babel}
\usepackage[letterpaper,margin=2cm,headheight=15pt,headsep=0.5cm,footskip=1cm]{geometry}
\usepackage{fancyhdr}
\usepackage{xcolor}

\pagestyle{fancy}
\fancyhf{}
\lhead{Rouleau-Desrochers, Christophe}
\lfoot{Reconnaissances}
\rfoot{\thepage}
\renewcommand{\headrulewidth}{0pt}

\setlength{\parindent}{0pt}
\setlength{\parskip}{2pt}

\newcommand{\groupe}[1]{\par\vspace{3pt}\textbf{#1}\par\vspace{-0.5pt}}

\begin{document}
\raggedright

\textbf{Reconnaissances}

\groupe{1. Bourses obtenues par concours fondé sur l'excellence}

\textbf{Bourse d'études supérieures en recherche du Canada au niveau de la maîtrise (BESRC~M)}, CRSNG, 27\,000~\$, 2025-2026. Concours national des trois organismes fédéraux. Les demandes sont évaluées par l'université d'accueil à l'intérieur d'un quota; la mienne l'a été par UBC, selon l'excellence universitaire (50~\%), le potentiel de recherche (30~\%) et les caractéristiques et habiletés interpersonnelles (20~\%). Une moyenne de première classe est exigée pour chacune des deux dernières années d'études.

\textbf{Bourses de recherche de 1\textsuperscript{er} cycle (BRPC)}, CRSNG, 6\,000~\$ par été, 2022, 2023 et 2024. Obtenues trois années de suite, au concours institutionnel de l'UQÀM en 2022, puis à celui de UBC en 2023 et 2024. Chaque université reçoit un quota de bourses du CRSNG et classe les candidatures selon l'excellence universitaire, le potentiel en recherche et la qualité attendue de la formation et du mentorat. À l'UQÀM, il faut une moyenne cumulative d'au moins 3,0 sur 4,3, et un comité de 15 à 25 chercheurs pondère ces critères à 40, 30 et 30~\%.

\textbf{Peter Rennie Memorial Award}, UBC, 1\,250~\$, 2025-2026. Bourse de la Faculté de foresterie réservée aux étudiants des cycles supérieurs dont la recherche porte sur les aspects environnementaux des sols forestiers et des arbres, attribuée sur recommandation de la Faculté.

\textbf{Braham G. Griffith Scholarship}, UBC, 500~\$, 2025-2026. Bourse d'études supérieures de la Faculté de foresterie, attribuée sur recommandation de la Faculté.

\textbf{Bourse Cocodet en sciences}, Fondation de l'UQÀM, 21\,375~\$ par année, 2022-2023 et 2023-2024. Obtenue au concours de la Fondation de l'UQÀM, sur dossier de candidature et sur la base de l'excellence scolaire.

\textbf{Supplément à l'appui de la diversité de l'effectif dans le secteur forestier canadien}, CRSNG et Service canadien des forêts, 5\,000~\$, été 2023. Ajouté à ma BRPC de 2023, sans demande distincte. Le supplément vise des étudiants de haut calibre : le CRSNG considère d'office les titulaires de BRPC ou de BESRC dont le projet touche un domaine d'intérêt pour le Service canadien des forêts et retient les projets les plus liés au secteur forestier. Le programme vise aussi à diversifier la main-d'œuvre forestière et peut prioriser les groupes sous-représentés. Il y en avait 30 en 2023 au Canada.

\textbf{Bourse de la Banque TD}, Fondation de l'UQÀM, 2\,000~\$, hiver 2023. Obtenue au concours de la Fondation de l'UQÀM, sur dossier de candidature et sur la base de l'excellence scolaire.

% Trottier : deux bourses de 1 500 $ ou doublon dans le portail?
\textbf{Bourse Jeanne-D'Arc et Bertin-Trottier en sciences}, Fondation de l'UQÀM, 1\,500~\$, 2022-2023. Bourse d'excellence destinée aux personnes qui reviennent aux études après quelques années sur le marché du travail, attribuée au concours de la Fondation. C'était mon cas, après presque dix ans en restauration.

\groupe{2. Distinctions attribuées d'office, sur la seule base des résultats scolaires}

\textbf{Médaille académique d'argent du Gouverneur général du Canada}, UQÀM, 2024. Remise à la personne finissante du baccalauréat qui a obtenu les meilleurs résultats. L'université la décerne sur le seul rendement scolaire, sans candidature, et le nombre de médailles dépend de sa taille : en 2024, l'UQÀM en a remis deux pour l'ensemble de ses finissants au baccalauréat. Obtenue au terme de mon B.Sc. en biologie; aucune somme n'y est rattachée.

\textbf{Liste d'excellence du doyen}, Faculté des sciences de l'UQÀM, 2021-2024. Deux fois par année, la Faculté y inscrit d'office les étudiants au baccalauréat classés dans les 5~\% supérieurs de leur programme.

\end{document}